% =========================================================================
% Chapter 1 - Kohn-Sham DFT
% =========================================================================
\chapter{Kohn--Sham density-functional theory}
\label{ch:ksdft}

This chapter fixes the physical objects the rest of the guide manipulates. The
emphasis is on \emph{which statements are theorems and which are
approximations}, because that distinction determines what a neural network may
legitimately be asked to learn, and what may be enforced on it as a hard
constraint.

\section{The many-body problem}

The non-relativistic electronic Hamiltonian for $N$ electrons in the field of
fixed nuclei is, in Hartree atomic units,
\begin{equation}
  \hat{H}
  = -\tfrac{1}{2}\sum_{i}^{N}\nabla_i^2
    + \sum_{i}^{N}\vext(\rr_i)
    + \sum_{i<j}\frac{1}{|\rr_i-\rr_j|} ,
  \label{eq:hamiltonian}
\end{equation}
where $\vext$ collects the electron--nucleus attraction. The wavefunction
$\Psi(\rr_1,\dots,\rr_N)$ depends on $3N$ coordinates, so a direct solution is
hopeless beyond a handful of electrons.

\section{The Hohenberg--Kohn theorems}
\label{sec:hk}

Density-functional theory replaces $\Psi$ by the electron density
\begin{equation}
  \rho(\rr) = N\int |\Psi(\rr,\rr_2,\dots,\rr_N)|^2\,
              \mathrm{d}\rr_2\cdots\mathrm{d}\rr_N ,
\end{equation}
a function of three coordinates. Two theorems justify this.

\begin{description}
  \item[First theorem.] The ground-state density determines the external
    potential up to an additive constant. Consequently, for a fixed electron
    number, the ground-state density is a functional of $\vext$ alone:
    \begin{equation}
      \vext(\rr) \;\longleftrightarrow\; \rho_0(\rr).
      \label{eq:hkmap}
    \end{equation}
  \item[Second theorem.] There exists a universal functional $F[\rho]$ such that
    the total energy
    \begin{equation}
      E[\rho] = F[\rho] + \int \vext(\rr)\rho(\rr)\,\mathrm{d}\rr
      \label{eq:energyfunctional}
    \end{equation}
    is minimised, subject to $\int\rho\,\mathrm{d}\rr = N$, by the ground-state
    density, where it equals the ground-state energy.
\end{description}

Equation~\eqref{eq:hkmap} is the formal justification for the first learned
model of this work. The target of a $\vext \mapsto \rho$ regression is not a
fitted correlation: it is a single-valued function whose existence and
uniqueness are guaranteed. That is a stronger footing than most machine-learning
targets enjoy.

\begin{pwarn}
The map is conditional on the electron number $N$. In the present data set
$N=\sum_a Z^{\mathrm{val}}_a$ follows from the geometry, so a network can infer
it; for charged systems or a chemically heterogeneous data set, $N$ must become
an explicit input.
\end{pwarn}

\section{The Kohn--Sham construction}

The difficulty is that $F[\rho]$ is unknown, and its dominant part --- the
kinetic energy --- is badly approximated by explicit density functionals. Kohn
and Sham sidestep this by introducing an auxiliary system of non-interacting
electrons with the \emph{same} density, for which the kinetic energy can be
computed exactly from orbitals. Writing
\begin{equation}
  F[\rho] = \Ts[\rho] + E_{\mathrm{H}}[\rho] + E_{\mathrm{xc}}[\rho],
\end{equation}
with the Hartree energy
\begin{equation}
  E_{\mathrm{H}}[\rho] = \frac{1}{2}\iint
    \frac{\rho(\rr)\rho(\rr')}{|\rr-\rr'|}\,\mathrm{d}\rr\,\mathrm{d}\rr',
\end{equation}
the variational condition becomes a set of single-particle equations,
\begin{equation}
  \left(-\tfrac{1}{2}\nabla^2 + v_{\mathrm{s}}[\rho](\rr)\right)\psi_i(\rr)
    = \varepsilon_i \psi_i(\rr),
  \qquad
  v_{\mathrm{s}} = \vext + \vH[\rho] + \vxc[\rho],
  \label{eq:kseq}
\end{equation}
with
\begin{equation}
  \rho(\rr) = \sum_i^{\mathrm{occ}} f_i |\psi_i(\rr)|^2,
  \qquad
  \Ts = -\tfrac{1}{2}\sum_i^{\mathrm{occ}} f_i
        \langle \psi_i | \nabla^2 | \psi_i \rangle .
\end{equation}
Because $v_{\mathrm{s}}$ depends on $\rho$, which depends on the orbitals,
Eq.~\eqref{eq:kseq} must be solved self-consistently.

\subsection{Why the map is hard}

Two features of Eq.~\eqref{eq:kseq} shape the choice of neural architecture in
Chapter~\ref{ch:fno}.

\begin{enumerate}
  \item \textbf{It is non-local.} The Hartree term is a
    $1/|\rr-\rr'|$ convolution: the density at one point depends on the
    potential everywhere. A network with a finite receptive field must be made
    very deep to propagate that information across a unit cell.
  \item \textbf{It is self-consistent.} The density is a fixed point, not a
    single forward evaluation. A supervised model learns the fixed point
    directly, which is what makes it fast --- and what makes it fragile off the
    training distribution.
\end{enumerate}

\section{The kinetic energy density}
\label{sec:tau}

The quantity stored in \file{TAUCAR} is the \emph{positive-definite} kinetic
energy density of the Kohn--Sham system,
\begin{equation}
  \tau(\rr) = \frac{1}{2}\sum_i^{\mathrm{occ}} f_i
              \bigl|\nabla\psi_i(\rr)\bigr|^2 .
  \label{eq:taudef}
\end{equation}
It integrates to $\Ts$ and is manifestly non-negative. An alternative
definition using the Laplacian,
$\tilde\tau = -\tfrac{1}{2}\sum_i f_i \psi_i^*\nabla^2\psi_i$, integrates to the
same $\Ts$ but differs pointwise by $\tfrac{1}{4}\nabla^2\rho$ and is
\emph{not} sign-definite. The distinction matters here: the constraints of
Section~\ref{sec:constraints} apply to Eq.~\eqref{eq:taudef}, and applying them
to the Laplacian form would be simply wrong.

\begin{ptip}
Chapter~\ref{ch:vasp} shows, from the VASP source, that \file{TAUCAR} indeed
stores Eq.~\eqref{eq:taudef} --- built spectrally as $i(\GG+\kk)$ acting on the
plane-wave coefficients --- and not the Laplacian form. This was verified
rather than assumed.
\end{ptip}

\section{Exact properties of \texorpdfstring{$\tau$}{tau}}
\label{sec:constraints}

Several relations constrain $\tau$ exactly. They are theorems, and
Chapter~\ref{ch:pifno} uses them as hard constraints rather than as penalties.

\begin{description}
  \item[Non-negativity.] $\tau(\rr)\ge 0$, immediately from
    Eq.~\eqref{eq:taudef}.
  \item[The von Weizsäcker bound.] For any system,
    \begin{equation}
      \tau(\rr) \;\ge\; \tauvW(\rr)
        \;\equiv\; \frac{|\nabla\rho(\rr)|^2}{8\rho(\rr)} ,
      \label{eq:hobound}
    \end{equation}
    with equality for a single nodeless orbital. This is the
    Hoffmann-Ostenhof inequality, and it is the single most useful constraint
    available for the $\rho\mapsto\tau$ map: $\tauvW$ is a closed-form
    functional of the input, so the bound can be imposed \emph{by
    construction}.
  \item[Uniform-gas limit.] As $\nabla\rho \to 0$,
    \begin{equation}
      \tau \;\longrightarrow\; \tauTF = C_{\mathrm{TF}}\,\rho^{5/3},
      \qquad
      C_{\mathrm{TF}} = \tfrac{3}{10}(3\pi^2)^{2/3} \approx 2.8712 .
    \end{equation}
  \item[Coordinate scaling.] For $\rho_\lambda(\rr)=\lambda^3\rho(\lambda\rr)$,
    \begin{equation}
      \Ts[\rho_\lambda] = \lambda^2\,\Ts[\rho] .
      \label{eq:scaling}
    \end{equation}
    This is exact and label-free: any training sample can be rescaled and its
    target $\Ts$ is then known in closed form, which is a data augmentation at
    zero DFT cost.
\end{description}

The decomposition suggested by Eq.~\eqref{eq:hobound},
\begin{equation}
  \tau = \tauvW[\rho] + \tauP,
  \qquad \tauP \ge 0 ,
  \label{eq:pauli}
\end{equation}
defines the \emph{Pauli term} $\tauP$. It is the part of the kinetic energy
that arises purely from Fermi statistics --- what distinguishes a many-fermion
system from a single-orbital one --- and it is the object every orbital-free
kinetic functional is really trying to approximate. Section~\ref{sec:pauli-head}
shows that learning $\tauP$ instead of $\tau$ is both more accurate and
cheaper.

\section{Pseudopotentials}

Core electrons are chemically inert but carry sharp oscillations that would
demand an impractically fine grid. A pseudopotential replaces the nucleus and
core by an effective ion of charge $Z^{\mathrm{val}}$ acting on smooth
pseudo-valence orbitals. Its local part behaves as
$-Z^{\mathrm{val}}e^2/r$ outside a core radius $R_{\mathrm{core}}$ and is
softened inside.

Three consequences run through this document:

\begin{enumerate}
  \item The fields in \file{CHGCAR} and \file{TAUCAR} are
    \emph{pseudo} quantities. They are mutually consistent, which is why the
    bound~\eqref{eq:hobound} is observed to hold on this data even though it is
    strictly guaranteed only for all-electron quantities.
  \item The external potential in \file{EXTCAR} is the \emph{local} part only.
    Non-local projectors are not functions of $\rr$ at all and cannot appear in
    a volumetric file even in principle.
  \item The short-range shape of the local pseudopotential is tabulated, not
    analytic. Chapter~\ref{ch:vasp} shows that modelling it with a smooth
    analytic form factor is the dominant error in a naive reconstruction, and
    how to avoid it.
\end{enumerate}
